\documentclass[11pt,a4paper]{article}

\usepackage[T1]{fontenc}
\usepackage[utf8]{inputenc}
\usepackage[english]{babel}
\usepackage{lmodern}
\usepackage{microtype}
\usepackage[a4paper,margin=20mm,headheight=15pt]{geometry}
\usepackage{graphicx}
\usepackage{booktabs}
\usepackage{tabularx}
\usepackage{array}
\usepackage{xcolor}
\usepackage{amsmath}
\usepackage{caption}
\usepackage{float}
\usepackage{enumitem}
\usepackage{fancyhdr}
\usepackage[hidelinks]{hyperref}
\usepackage{xurl}

\graphicspath{{figures/}}
\definecolor{controlblue}{HTML}{2563EB}
\definecolor{acceptedgreen}{HTML}{16A34A}
\definecolor{rejectedred}{HTML}{DC2626}
\definecolor{amber}{HTML}{D97706}
\definecolor{slate}{HTML}{334155}
\definecolor{lightgray}{HTML}{F1F5F9}

\newcommand{\dstwr}{DS-TWR}
\newcommand{\code}[1]{\texttt{\detokenize{#1}}}
\newcolumntype{Y}{>{\raggedright\arraybackslash}X}

\pagestyle{fancy}
\fancyhf{}
\lhead{\small Native DS-TWR round-boundary investigation}
\rhead{\small July 31, 2026}
\cfoot{\thepage}

\hypersetup{
  pdftitle={Native DS-TWR round-boundary investigation},
  pdfauthor={UWB ESP-IDF technical report},
  pdfsubject={Root cause and guarded validation at the 53 ms profile},
  pdfkeywords={UWB, DS-TWR, DW3000, ESP32-S3, round boundary, timing integrity}
}

\begin{document}
\pagenumbering{gobble}

\begin{titlepage}
  \centering
  \vspace*{14mm}
  {\Huge\bfseries Native \dstwr{} round-boundary investigation\par}
  \vspace{4mm}
  {\Large Why shorter frames lost integrity, and why the guarded 53 ms profile works\par}
  \vspace{14mm}

  \begin{tabular}{rl}
    Test date: & July 31, 2026, EEST\\
    Platform: & 5 ESP32-S3 + DW3000 modules, Raspberry Pi dashboard\\
    Branch: & \texttt{experiment/native-ds-round-boundary}\\
    Reference: & \texttt{c727a64} before boundary instrumentation\\
    Configuration: & M1 active tag; M2, M3, M4, M5 anchors\\
    Geometry: & unchanged static indoor setup\\
    Captures: & five uninterrupted one-hour blocks\\
  \end{tabular}

  \vfill
  \colorbox{lightgray}{%
    \begin{minipage}{0.88\textwidth}
      \vspace{3mm}
      \textbf{Key result.}
      The apparent speed/precision limit was not the airtime of an ordinary
      \texttt{POLL--RESP--FINAL} exchange. All \textbf{267} impossible ToFs
      observed before the fix occurred at the separately inserted
      anchor-geometry boundary; \textbf{none} occurred in an ordinary tag
      ranging round. A 5 ms guard on each side of that periodic exchange
      eliminated every source rejection during a new one-hour run at the
      fastest tested \textbf{53 ms} profile. Delivered throughput changed
      only from \textbf{67.24} to \textbf{65.35 ranges/s} ($-2.8\%$), while
      RESP timeouts fell by \textbf{53.7\%} and FINAL timeouts by
      \textbf{96.6\%}.
      \vspace{3mm}
    \end{minipage}}
  \vfill
  {\large Reproducible field investigation and validation report\par}
  \vspace{5mm}
\end{titlepage}

\pagenumbering{roman}
\tableofcontents
\clearpage
\pagenumbering{arabic}

\section{Executive summary}

The original hypothesis was that shorter Native \dstwr{} profiles reduced
accuracy because an individual radio exchange no longer fit inside its slot.
The instrumentation added for this study shows a more specific failure:
ordinary tag exchanges were sound, but the periodic transition into and out
of anchor self-localization was not isolated from adjacent radio work.

Five one-hour captures were collected. Four unguarded profiles varied the
normal frame from 59 ms down to 53 ms. The most aggressive 53 ms profile was
then repeated after moving the normal gap directly behind the tag round and
adding a 5 ms guard before and after the geometry exchange.

\begin{table}[H]
  \centering
  \caption{Five-hour result summary. Rates and event counts come directly
  from the lossless firmware log stream.}
  \small
  \begin{tabular}{lrrrrr}
    \toprule
    Normal frame & Valid ranges & Rate & RESP timeout & FINAL timeout & Source reject\\
    \midrule
    59 ms & 220,710 & 61.31/s & 4,657 & 350 & 34\\
    58 ms & 223,776 & 62.16/s & 4,972 & 376 & 35\\
    57 ms & 226,315 & 62.86/s & 6,660 & 858 & 98\\
    53 ms & 242,069 & 67.24/s & 8,514 & 923 & 100\\
    \textcolor{acceptedgreen}{\textbf{53 ms guarded}}
      & \textbf{235,249} & \textbf{65.35/s} & \textbf{3,944}
      & \textbf{31} & \textbf{0}\\
    \bottomrule
  \end{tabular}
\end{table}

\paragraph{Decision.}
The guarded 53 ms profile is a valid fast operating point for the tested
static setup. It preserves the short ordinary frame while spending extra time
only around the periodic geometry-maintenance operation. The validation is
strong evidence of timing integrity, but not yet a universal RF or motion
qualification; dynamic, NLOS, and longer-duration trials should follow.

\section{Question under test}

The Native \dstwr{} tag exchanges three mandatory packets with each anchor:
\texttt{POLL}, \texttt{RESP}, and \texttt{FINAL}. With four anchors, one
normal round contains four independent three-message exchanges. A separate
anchor-to-anchor \dstwr{} operation is periodically inserted so the dashboard
can reconstruct a moving physical anchor geometry.

The investigation separated three possible causes:

\begin{enumerate}[leftmargin=8mm]
  \item \textbf{insufficient ordinary slot duration}: the packet exchange
  itself overruns its allocation;
  \item \textbf{anchor-specific behavior}: A2, its RF path, or its antenna
  calibration is worse than the other anchors;
  \item \textbf{round-boundary state overlap}: timestamps or radio state from
  the tag round and the separately inserted geometry exchange touch each
  other at the schedule boundary.
\end{enumerate}

The third hypothesis is the only one consistent with all captured evidence.

\section{Instrumentation and integrity checks}

The firmware was instrumented at the radio boundary rather than only in the
dashboard solver. Every exchange carries and validates:

\begin{itemize}[leftmargin=7mm]
  \item random token, 16-bit sequence, round, slot, and first-anchor ID;
  \item flags identifying a geometry exchange and the first following tag
  round;
  \item DW3000 transmit/receive timestamps for all stages;
  \item low-32 ESP timer values for host-side scheduling correlation;
  \item separate alternating initiator and responder exchange contexts.
\end{itemize}

Before applying the \dstwr{} equation, the code requires exact context
equality and valid timestamp intervals. A non-finite, non-positive, negative,
or physically impossible ToF is rejected at the source before telemetry and
position solving. Consequently, a clean trail cannot hide an invalid raw
measurement behind later solver rejection.

The first anchor also rotates between rounds. If the failure remained tied to
A2, that would indicate an anchor path. If it followed schedule position and
appeared on all anchors, it would indicate a boundary problem.

\section{Results before the guard}

\subsection{Throughput}

Shortening the frame increases delivered ranges as expected. The pre-guard
53 ms profile reaches 67.24 valid ranges/s; the guarded validation reaches
65.35 ranges/s. The small difference is the cost of periodic boundary
isolation, not an increase in the ordinary tag frame.

\begin{figure}[H]
  \centering
  \includegraphics[width=0.93\textwidth]{01_valid_throughput.pdf}
  \caption{Delivered raw-range throughput. Blue denotes the original
  schedule; green denotes the corrected 53 ms validation.}
\end{figure}

\subsection{Every impossible ToF is a boundary event}

The decisive observation is categorical rather than statistical. Across the
four original one-hour runs, all 267 source rejects occurred either inside
the geometry exchange or in the first tag round immediately after it. There
were zero source rejects in ordinary rounds.

\begin{table}[H]
  \centering
  \caption{Location of every rejected ToF before and after the guard.}
  \small
  \begin{tabular}{lrrrr}
    \toprule
    Frame & Geometry exchange & First post-geometry round & Ordinary round & Total\\
    \midrule
    59 ms & 29 & 5 & 0 & 34\\
    58 ms & 25 & 10 & 0 & 35\\
    57 ms & 81 & 17 & 0 & 98\\
    53 ms & 78 & 22 & 0 & 100\\
    53 ms guarded & 0 & 0 & 0 & 0\\
    \bottomrule
  \end{tabular}
\end{table}

\begin{figure}[H]
  \centering
  \includegraphics[width=0.93\textwidth]{02_boundary_rejects.pdf}
  \caption{Source-rejected negative ToFs classified by scheduler location.
  The absence of an ordinary-round segment is the key result.}
\end{figure}

\subsection{The problem is not A2}

The rejected measurements appear on A2, A3, A4, and A5. Their distribution
changes when the first-anchor order rotates. A2 has the fewest aggregated
rejects, while A5 has the most; the fault therefore follows timing position,
not one permanent anchor identity.

\begin{figure}[H]
  \centering
  \includegraphics[width=0.93\textwidth]{05_anchor_distribution.pdf}
  \caption{Rejected measurements by responding anchor. Rotation moves the
  boundary exposure across the four modules and disproves an A2-only fault.}
\end{figure}

\section{Root cause}

In the original scheduler, the configured round gap was applied after the
optional geometry operation. It therefore did not isolate the last tag
\texttt{FINAL} from the following \texttt{SURVEY\_CMD}. On the other side of
the geometry slot, the next tag \texttt{POLL} relied only on the nominal slot
duration. Occasional radio and host scheduling latency allowed the two
initiator regimes to touch at either boundary.

This explains why packet-airtime arithmetic alone was insufficient. An
ordinary three-message exchange could fit comfortably while the separately
inserted maintenance exchange still collided with the state transition around
it. The resulting failure was rare, short, and often invisible after solver
gating, but it was present in the raw timestamps.

\begin{figure}[H]
  \centering
  \includegraphics[width=0.96\textwidth]{06_schedule_guard.pdf}
  \caption{Conceptual scheduler comparison. The correction protects both
  sides of geometry maintenance without lengthening an ordinary 53 ms tag
  round.}
\end{figure}

\section{Correction}

The corrected schedule makes five changes:

\begin{enumerate}[leftmargin=8mm]
  \item place the configured frame gap immediately after every tag round;
  \item wait 5 ms before transmitting \texttt{SURVEY\_CMD};
  \item wait another 5 ms after the anchor-geometry \dstwr{} slot;
  \item mark every slot in the first following tag round as post-geometry;
  \item retain complete diagnostic traces for every rejected exchange while
  sampling detailed accepted traces to avoid perturbing the scheduler.
\end{enumerate}

The collector was also corrected to deduplicate by token, round, and slot
rather than by the wrapping 16-bit sequence alone. The final one-hour run
crossed several sequence wraps while its count remained linear.

\section{Guarded 53 ms validation}

The final firmware was rebuilt, uploaded simultaneously to all five modules,
and tested for one uninterrupted hour at slot 13 ms plus gap 1 ms. It produced
235,249 accepted ranges, no impossible ToF, no timestamp-context mismatch,
and no ordinary-round integrity error.

\begin{figure}[H]
  \centering
  \includegraphics[width=0.93\textwidth]{04_before_after_timeline.pdf}
  \caption{Per-minute delivery and cumulative source rejects for the two 53
  ms runs. The guarded capture remains at zero rejects for the full hour.}
\end{figure}

The guard reduces valid throughput by 2.8\%, from 67.24 to 65.35 ranges/s.
That cost is smaller than the nominal added guard time would suggest because
fewer collided exchanges are retried or lost. RESP timeouts fall from 8,514
to 3,944 ($-53.7\%$), FINAL timeouts fall from 923 to 31 ($-96.6\%$), and
source-rejected ToFs fall from 100 to zero.

\begin{figure}[H]
  \centering
  \includegraphics[width=0.93\textwidth]{03_failure_rates.pdf}
  \caption{Failure events normalized per 1,000 delivered ranges. The guarded
  schedule improves both integrity and radio-exchange completion.}
\end{figure}

\section{How to interpret the result}

\subsection{What is demonstrated}

\begin{itemize}[leftmargin=7mm]
  \item The ordinary 53 ms Native \dstwr{} frame is not the source of the
  observed impossible ranges in this static test.
  \item The prior apparent accuracy limit was a periodic scheduler-boundary
  problem, not a fundamental \dstwr{} propagation or airtime limit.
  \item The error is not specific to A2, antenna delay, or one RF path.
  \item Source validation prevents a rare timestamp corruption from becoming
  a plausible-looking distance or position.
  \item The boundary guards buy clean integrity for only 2.8\% measured
  delivered-throughput cost.
\end{itemize}

\subsection{What is not yet demonstrated}

This test does not establish a universal 53 ms limit. It uses one static
indoor geometry and one RF environment. A clean source-reject count proves
timestamp and schedule integrity, but it is not by itself a position-accuracy
metric. The next qualification should include moving-tag trials, controlled
NLOS conditions, power/temperature variation, and multiple runs of at least
one hour. Position RMSE and P95 should be compared against the same fixed
geometric reference used for the slow control.

\subsection{Recommended operating policy}

Use the guarded 53 ms profile as the current fast candidate. Keep the 59 ms
and 58 ms profiles available as diagnostic controls, but the present evidence
does not justify slowing ordinary tag rounds merely to hide a geometry
boundary that can be isolated directly. Preserve all source context checks
and physically impossible-ToF rejection in production builds.

\section{Reproducibility and artifacts}

The complete raw captures, metadata snapshots, generated figures, and
machine-readable report summary live under:

\begin{quote}
  \code{reports/native_ds_round_boundary_20260731/}
\end{quote}

The classification can be reproduced with:

\begin{verbatim}
tools/analyze_native_ds_round_boundary.py \
  reports/native_ds_round_boundary_20260731/data/\
candidate_59ms_round_boundary_3600s.logs.jsonl \
  reports/native_ds_round_boundary_20260731/data/\
candidate_58ms_round_boundary_3600s.logs.jsonl \
  reports/native_ds_round_boundary_20260731/data/\
candidate_57ms_round_boundary_3600s.logs.jsonl \
  reports/native_ds_round_boundary_20260731/data/\
candidate_53ms_round_boundary_3600s.logs.jsonl \
  reports/native_ds_round_boundary_20260731/data/\
candidate_53ms_guarded_3600s.logs.jsonl
\end{verbatim}

Figures and \code{report_summary.json} are regenerated by:

\begin{verbatim}
tools/uwb_native_ds_round_boundary_report.py
\end{verbatim}

The interrupted 130 s over-instrumented capture is retained separately and is
not included in any result above.

\end{document}
